\KOMAoption{paper}{a4, portrait}
\areaset{267mm}{180mm}
\fancyheadoffset{0pt} % обновляем хедер и футер

\newgeometry{
  top=14mm,
  bottom=11mm,
  right=10mm,
  left=22mm,
}

\begin{landscape}
\setlength{\footskip}{15.5pt}

\color{unidarkgreen}{\section[(обязательное) Схема электрическая принципиальная]{(обязательное)\\Схема электрическая принципиальная}\label{app:a3}}
\color{black}

\setlength{\footskip}{15.5pt}
\vskip 3mm

\begin{figure}[H]
  \centering
  \includegraphics[width=1.1\textwidth,height=1.1\textheight,keepaspectratio,trim=5pt 5pt 5pt 5pt,clip]{\apppath/Приложение. Схема Э3/_latex/img/Fig_1.pdf}
  \captionof{figure}{Схема электрическая принципиальная (часть 1)}\label{appvo:fig_1}
\end{figure}

\begin{figure}[H]
  \centering
  \includegraphics[width=1.3\textwidth,height=1.3\textheight,keepaspectratio,trim=5pt 5pt 5pt 5pt,clip]{\apppath/Приложение. Схема Э3/_latex/img/Fig_2.pdf}
  \captionof{figure}{Схема электрическая принципиальная (часть 2)}\label{appvo:fig_2}
\end{figure}

\begin{figure}[H]
  \centering
  \includegraphics[width=1.3\textwidth,height=1.3\textheight,keepaspectratio,trim=5pt 5pt 5pt 5pt,clip]{\apppath/Приложение. Схема Э3/_latex/img/Fig_3.pdf}
  \captionof{figure}{Схема электрическая принципиальная (часть 3)}\label{appvo:fig_3}
\end{figure}

\begin{figure}[H]
  \centering
  \includegraphics[width=1.3\textwidth,height=1.3\textheight,keepaspectratio,trim=5pt 5pt 5pt 5pt,clip]{\apppath/Приложение. Схема Э3/_latex/img/Fig_4.pdf}
  \captionof{figure}{Схема электрическая принципиальная (часть 4)}\label{appvo:fig_4}
\end{figure}

\begin{figure}[H]
  \centering
  \includegraphics[width=1.3\textwidth,height=1.3\textheight,keepaspectratio,trim=5pt 5pt 5pt 5pt,clip]{\apppath/Приложение. Схема Э3/_latex/img/Fig_5.pdf}
  \captionof{figure}{Схема электрическая принципиальная (часть 5)}\label{appvo:fig_5}
\end{figure}

\begin{figure}[H]
  \centering
  \includegraphics[width=1.3\textwidth,height=1.3\textheight,keepaspectratio,trim=5pt 5pt 5pt 5pt,clip]{\apppath/Приложение. Схема Э3/_latex/img/Fig_6.pdf}
  \captionof{figure}{Схема электрическая принципиальная (часть 6)}\label{appvo:fig_6}
\end{figure}

\begin{figure}[H]
  \centering
  \includegraphics[width=1.3\textwidth,height=1.3\textheight,keepaspectratio,trim=5pt 5pt 5pt 5pt,clip]{\apppath/Приложение. Схема Э3/_latex/img/Fig_7.pdf}
  \captionof{figure}{Схема электрическая принципиальная (часть 7)}\label{appvo:fig_7}
\end{figure}

\begin{figure}[H]
  \centering
  \includegraphics[width=1.3\textwidth,height=1.3\textheight,keepaspectratio,trim=5pt 5pt 5pt 5pt,clip]{\apppath/Приложение. Схема Э3/_latex/img/Fig_8.pdf}
  \captionof{figure}{Схема электрическая принципиальная (часть 8)}\label{appvo:fig_8}
\end{figure}

\begin{figure}[H]
  \centering
  \includegraphics[width=1.3\textwidth,height=1.3\textheight,keepaspectratio,trim=5pt 5pt 5pt 5pt,clip]{\apppath/Приложение. Схема Э3/_latex/img/Fig_9.pdf}
  \captionof{figure}{Схема электрическая принципиальная (часть 9)}\label{appvo:fig_9}
\end{figure}

\begin{figure}[H]
  \centering
  \includegraphics[width=1.3\textwidth,height=1.3\textheight,keepaspectratio,trim=5pt 5pt 5pt 5pt,clip]{\apppath/Приложение. Схема Э3/_latex/img/Fig_10.pdf}
  \captionof{figure}{Схема электрическая принципиальная (часть 10)}\label{appvo:fig_10}
\end{figure}

\begin{figure}[H]
  \centering
  \includegraphics[width=1.3\textwidth,height=1.3\textheight,keepaspectratio,trim=5pt 5pt 5pt 5pt,clip]{\apppath/Приложение. Схема Э3/_latex/img/Fig_11.pdf}
  \captionof{figure}{Схема электрическая принципиальная (часть 11)}\label{appvo:fig_11}
\end{figure}
\clearpage

\end{landscape}
\restoregeometry